% CV of Baye Lahad Mbacke -- build with: pdflatex Resume_PHD_en.tex
\documentclass[11pt,a4paper]{article}
\usepackage[utf8]{inputenc}
\usepackage[T1]{fontenc}
\usepackage{lmodern}
\usepackage[margin=1.8cm,top=1.5cm,bottom=1.5cm]{geometry}
\usepackage{xcolor}
\usepackage{tabularx}
\usepackage[hidelinks,pdftitle={Baye Lahad Mbacke -- CV},pdfauthor={Baye Lahad Mbacke}]{hyperref}

\definecolor{accent}{HTML}{0066CC}
\pagestyle{empty}
\setlength{\parindent}{0pt}

\newcommand{\cvsection}[1]{%
  \vspace{10pt}{\Large\color{accent}#1}\par\vspace{-6pt}%
  {\color{accent}\rule{\linewidth}{0.8pt}}\par\vspace{4pt}}

% \cventry{institution}{location}{title}{dates}
\newcommand{\cventry}[4]{%
  \begin{tabularx}{\linewidth}{@{}X r@{}}
    \textbf{#1} & #2 \\
    \textit{#3} & \textit{#4}
  \end{tabularx}\par}

\newenvironment{cvitems}{%
  \begin{list}{--}{\setlength{\leftmargin}{1.5em}\setlength{\itemsep}{1pt}
  \setlength{\parsep}{0pt}\setlength{\topsep}{2pt}}}{\end{list}\vspace{4pt}}

\begin{document}

\begin{center}
  {\fontsize{24}{28}\selectfont\color{accent} Baye Lahad Mbacke}\\[6pt]
  PhD Student in Computer Science -- Nantes University, LS2N\\[2pt]
  8 boulevard Mendès France, Nantes, France \textbullet{} +33 6 01 32 80 22\\[2pt]
  \href{mailto:lahadmbacke00@gmail.com}{lahadmbacke00@gmail.com} \textbullet{}
  \href{https://lahadmbacke.github.io/Lahad_Mbacke/}{Portfolio} \textbullet{}
  \href{https://linkedin.com/in/baye-lahad-mbacke}{LinkedIn} \textbullet{}
  \href{https://github.com/LahadMbacke}{GitHub} \textbullet{}
  \href{https://huggingface.co/Lahad}{Hugging Face}
\end{center}

\vspace{4pt}
PhD student at Nantes University since March 2026, within the Laboratoire des Sciences du
Numérique de Nantes (LS2N). My research focuses on efficient and frugal methods for the automatic
extraction and structuring of cultural event information, in order to connect cultural and
creative industries with their audiences.

\cvsection{Education}

\cventry{Nantes University -- LS2N (TALN team) \& LIA}{Nantes, France}
        {PhD in Computer Science (MaSTIC Doctoral School)}{Mar. 2026--Present}
\textit{Automatic extraction and structuring of cultural events: an efficient and frugal approach
to connect cultural and creative industries with their audiences}
\begin{cvitems}
  \item Supervisors: Richard Dufour (LS2N, Nantes University) and Vincent Labatut (LIA, Avignon University)
  \item Industrial partner: ideactiv (French cultural indexing platform)
  \item Detecting event listings on cultural websites, identifying individual events, extracting
        and consolidating their information (title, date, venue, price, booking, etc.) without
        hallucinations or costly third-party services
\end{cvitems}

\cventry{Nantes University}{Nantes, France}
        {Master ATAL (Machine Learning and Natural Language Processing)}{2023--2025}
Specialization in advanced NLP techniques and machine learning\par\vspace{6pt}

\cventry{University of Strasbourg}{Strasbourg, France}{Bachelor's Degree in Computer Science (L3)}{2022--2023}
Advanced study in algorithms and programming\par\vspace{6pt}

\cventry{Cheikh Anta Diop University of Dakar}{Dakar, Senegal}
        {Bachelor's Degree in Mathematics and Computer Science}{2018--2022}
Fundamentals of mathematics and computer science

\cvsection{Professional Experience}

\cventry{Ouest-France}{Rennes, France}
        {Assistant Data Scientist -- Research \& Development (work-study)}{Sept. 2024--Sept. 2025}
\begin{cvitems}
  \item Geotagging: identifying location references from address fragments in news articles,
        and linking them to place tags
  \item Knowledge extraction: named entity recognition and entity relation mapping
  \item Disambiguation of geographical entities (addresses, squares, monuments, cities, etc.)
\end{cvitems}

\cventry{Laboratoire des Sciences du Numérique de Nantes (LS2N)}{Nantes, France}
        {Research Intern}{Jan. 2024--May 2024}
\begin{cvitems}
  \item Comparison of named entity recognition systems on press articles
  \item Machine learning for key information extraction
  \item Annotation and analysis of structured and unstructured data
\end{cvitems}

\cvsection{Academic Projects}

\cventry{Nantes University}{Nantes, France}
        {DEFT2013 Challenge -- Recipe classification}{Feb. 2024--Mar. 2024}
\begin{cvitems}
  \item Classification models for recipes (main course, starter, dessert) based on titles and
        instructions, using NLP and machine learning methods
\end{cvitems}

\cventry{University of Strasbourg}{Strasbourg, France}
        {Autoencoders and classification}{Apr. 2023--May 2023}
\begin{cvitems}
  \item Evaluation of autoencoders for tabular data classification
  \item Comparative analysis of performance with and without compression into a reduced latent space
\end{cvitems}

\cvsection{Technical Skills}
\begin{tabularx}{\linewidth}{@{}l X@{}}
  \textbf{Machine \& Deep Learning} & PyTorch, Scikit-learn, LLMs, Fine-tuning, Quantization, Distillation \\
  \textbf{NLP} & spaCy, NLTK, LangChain, RAG, Agents, Named Entity Recognition \\
  \textbf{Programming} & Python, R, Java, C/C++ \\
  \textbf{Databases} & SQL, MySQL, Oracle, PostgreSQL, MongoDB
\end{tabularx}

\cvsection{Languages}
\textbf{Wolof}: native \quad\textbullet\quad
\textbf{French}: fluent \quad\textbullet\quad
\textbf{English}: B2

\cvsection{Interests}
\textbf{Reading}: newspapers, tech news, technology watch \quad\textbullet\quad
\textbf{Sport}: football

\end{document}
